% STATUT : À RÉDIGER, mais TOUT est déjà calculé (scripts vasicek_lab).
% C'est le chapitre qui établit ce que la donnée FIXE, par opposition au
% chapitre identifiabilité qui établit ce qu'elle ne fixe pas.
\chapter{Données et ancrage empirique}
\label{ch:donnees}

\textit{[Chapitre à rédiger. Les résultats existent, il s'agit de les mettre en forme.]}

% PLAN :
%  1. Les sources : Data Breach Chronology (notifications), SAS OpRisk
%     Global Data (pertes opérationnelles Bâle), VCDB (sous-piliers).
%     Périmètre, limites, biais de collecte à annoncer d'emblée.
%  2. FRÉQUENCE d'entrée : lambda ~ 21,6/an, mélange de Poisson,
%     sur-dispersion phi ~ 9,2 (binomiale négative).
%  3. ACCUMULATION : concentration journalière, Gini 0,69 ; MOVEit comme
%     cause commune à 191 victimes. Poisson composé / structure de paquet.
%  4. POURQUOI PAS DE HAWKES : noyau exponentiel puis noyau à retard
%     (variante Bessy-Roland), excitation intra-journalière, endogénéité
%     qui tombe à zéro quand on isole l'exogène (logique two-phase).
%     Rejet ARGUMENTÉ, pas par ignorance : c'est un point fort.
%  5. SÉVÉRITÉ : POT/GPD, xi ~ 0,90, sensibilité au seuil (0,93 / 0,83 /
%     0,68 aux seuils q90, q95, q99), hétérogénéité par catégorie
%     (0,92 à 1,37) qui justifie un xi_j par pilier.
%  6. Sous-piliers (VCDB) : P2 et P3 atteignables sur open data,
%     P4 s'effondre (56 observations en finance). Conditionne le cas
%     d'usage approfondi aux données client.
%  7. Synthèse : le tableau de ce qui est calibré, avec son incertitude.
